% This document is compiled using pdfLaTeX
% Purpose: 给 Codex / Claude Code 等“读着写代码”的说明稿（对应论文第4章：加速方法）
% Usage: % This document is compiled using pdfLaTeX
% Purpose: 给 Codex / Claude Code 等“读着写代码”的说明稿（对应论文第4章：加速方法）
% Usage: % This document is compiled using pdfLaTeX
% Purpose: 给 Codex / Claude Code 等“读着写代码”的说明稿（对应论文第4章：加速方法）
% Usage: % This document is compiled using pdfLaTeX
% Purpose: 给 Codex / Claude Code 等“读着写代码”的说明稿（对应论文第4章：加速方法）
% Usage: \input{speed.tex} 或单独编译均可

\documentclass{article}
\usepackage{amsmath,amssymb,amsfonts,amsthm}
\usepackage{enumitem}
\usepackage{geometry}
\usepackage{mathrsfs}
\geometry{a4paper, margin=1in}
\usepackage{ctex} % 中文
\usepackage{hyperref}
\usepackage{float}
\usepackage[linesnumbered,ruled,vlined]{algorithm2e}

\newtheorem{remark}{注}
\newtheorem{definition}{定义}

\begin{document}

% ============================================================
%  Speed-up strategies for LBBD (Chapter 4)
%  1) Valid inequalities for BMP (master)
%  2) Warm-start initial cuts via dual-feasible solutions
% ============================================================

\section{LBBD 算法加速策略（第4章：有效不等式 + 初始割）}

\subsection{符号速查（只列本文件会用到的）}

\begin{itemize}[leftmargin=2em]
  \item 规划期：$T=\{1,\dots,l\}$。
  \item 收集点集合：$V_S$。在路由子问题中，客户集合写作 $N_c$（通常可取 $N_c=V_S$ 或“当期被访问的收集点”）。
  \item 车辆：可用车辆数上限 $K$，单车容量 $Q$。
  \item 主问题变量：
    \begin{itemize}[leftmargin=2em]
      \item $z_{it}\in\{0,1\}$：时期 $t$ 是否访问收集点 $i$。
      \item $\lambda_{ivt}\in\{0,1\}$：收集点 $i$ 的“上次访问时期为 $v$、本次访问时期为 $t$”的弧是否被选中（时间扩展网络/最短路式改写中的弧选择变量）。
      \item $\omega_t\ge 0$：时期 $t$ 的最优路径成本（由子问题给出，用 Benders 割约束它）。
    \end{itemize}
  \item 预计算参数（来自第2.4节改写模型；加速策略只\textbf{引用}它们，不重复推导）：
    \begin{itemize}[leftmargin=2em]
      \item $\pi(i,t)$：在“本次访问 $t$”条件下，\emph{可行的}最早上次访问时期下界（由收集点容量等逻辑得到）。
      \item $g_{ivt}$：若选择弧 $(v,t)$，即 $i$ 上次在 $v$ 被访问、当前在 $t$ 被访问，则时期 $t$ 从 $i$ 需要收走的\textbf{全量}（累计）收集量。
      \item （若实现 $T2$）$e_{ivt}$：与弧 $(v,t)$ 关联的收集点库存持有成本（或其它在第2.4节中被“弧化”的成本项）。
    \end{itemize}
\end{itemize}

\begin{remark}
本文件只覆盖第4章“加速方法”，默认你已经按 \texttt{lbbd\_all.tex} 中的说明实现了：
(1) BMP（主问题）中用 $\omega_t$ 代替显式路由变量；(2) BSP（子问题）用集合划分/列生成求解 CVRP(t) 并生成 Benders 割。
\end{remark}

% ------------------------------------------------------------
% A. Valid inequalities (add to BMP)
% ------------------------------------------------------------
\subsection{A. 主问题（BMP）有效不等式：直接加到主问题即可}

\subsubsection*{A.1 最小收集启动时间约束（式(4-1)--(4-2)）}

定义一个最早必须启动收集或生产的时期 $t'$：
\begin{equation}
t'=\arg\min_{1\le t\le l}\left\{\,P_{00}+\frac{I_{00}}{k}<\sum_{\tau=1}^{t} d_{\tau}\right\}.
\label{eq:tprime_def_41}
\end{equation}

若这样的 $t'$ 存在，则到 $t'$ 期末必须至少发生一次收集活动（否则没有新增原料来源，无法满足累积需求）：
\begin{equation}
\sum_{\tau=1}^{t'}\ \sum_{i\in V_S} z_{i\tau}\ \ge\ 1.
\label{eq:vi_start_collect_42}
\end{equation}

\paragraph{实现提示}
\begin{itemize}[leftmargin=2em]
  \item 这是一个\textbf{一次性}约束：只需在建模时计算 $t'$ 并加入一次 \eqref{eq:vi_start_collect_42}。
  \item 若不存在 $t'$（即初始库存 $P_{00}+I_{00}/k$ 足以覆盖全周期累积需求），则不要添加该约束。
\end{itemize}

\subsubsection*{A.2 周期车辆总容量约束（式(4-3)）}

时期 $t$ 的总计划收集量可由 $\lambda$ 的选择唯一确定（全量收集）：
\[
\text{Pickup}_t \;=\; \sum_{i\in V_S}\ \sum_{v=\pi(i,t)}^{t-1} g_{ivt}\,\lambda_{ivt}.
\]
由于每期最多使用 $K$ 辆容量为 $Q$ 的车，必要条件为：
\begin{equation}
\sum_{i\in V_S}\ \sum_{v=\pi(i,t)}^{t-1} g_{ivt}\,\lambda_{ivt}\ \le\ K\cdot Q,\qquad \forall t\in T.
\label{eq:vi_totalcap_43}
\end{equation}

\paragraph{实现提示}
\begin{itemize}[leftmargin=2em]
  \item 这是\textbf{按期}约束：对所有 $t\in T$ 添加 \eqref{eq:vi_totalcap_43}。
  \item 该不等式是可行性的\textbf{必要条件}（不是充分条件），但能显著收紧主问题的 LP 松弛。
\end{itemize}

\subsubsection*{A.3 大收集量客户约束（式(4-4)）}

若某个收集点在时期 $t$ 的收集量超过半车容量（$g_{ivt}>Q/2$），则同一条路径不可能再服务另一个同样“大”的收集点。
因此，在时期 $t$，满足 $g_{ivt}>Q/2$ 的被选中弧数量不能超过 $K$：
\begin{equation}
\sum_{i\in V_S}\ \sum_{\substack{v=\pi(i,t)\\ g_{ivt}>Q/2}}^{t-1}\ \lambda_{ivt}\ \le\ K,\qquad \forall t\in T.
\label{eq:vi_bigitem_44}
\end{equation}

\paragraph{实现提示}
\begin{itemize}[leftmargin=2em]
  \item 这同样是\textbf{按期}约束：对每个 $t$ 建一个线性约束，内部只枚举满足 $g_{ivt}>Q/2$ 的 $(i,v)$。
  \item 若存在 $g_{ivt}>Q$，则该弧本身会导致 CVRP(t) 不可行，建议在预处理阶段直接\textbf{删除}该弧或在主问题中禁用它（比加 \eqref{eq:vi_bigitem_44} 更“直接”）。
\end{itemize}

% ------------------------------------------------------------
% B. Initial cuts (warm start)
% ------------------------------------------------------------
\subsection{B. 主问题（BMP）初始割：Benders 迭代前的“预热”}

动机：LBBD 初期，$\omega_t$ 在主问题里可能被压到非常小（甚至 0），导致第一次/前几次迭代的主问题解偏离真实可行解，从而迭代次数增多。初始割为 $\omega_t$ 提供一个\textbf{有效下界}，加速收敛。

\subsubsection*{B.1 访问决策未知的 CVRP(t) 集合划分模型（式(4-5)）}

设 $\hat{\mathcal R}_t$ 为时期 $t$ 的所有可行路径集合（通过列生成隐式表示），路径 $r$ 的成本为 $b_r$，其对“客户-需求情景”$(i,v)$ 的覆盖系数为 $a_{ivr}\in\{0,1\}$。
则 CVRP(t) 可写为：
\begin{equation}
\begin{aligned}
\text{CVRP}(t)\;\; \omega_t=\min\ &\sum_{r\in\hat{\mathcal R}_t} b_r\,\xi_r\\
\text{s.t.}\ &\sum_{r\in\hat{\mathcal R}_t} a_{ivr}\,\xi_r=\lambda_{ivt},\quad \forall i\in N_c,\ v=\pi(i,t),\dots,t-1,\\
&\sum_{r\in\hat{\mathcal R}_t}\xi_r\le K,\\
&\xi_r\in\{0,1\},\quad \forall r\in\hat{\mathcal R}_t.
\end{aligned}
\label{eq:cvrp_sp_45}
\end{equation}

其 LP 松弛对偶为（式(4-6)）：
\begin{equation}
\begin{aligned}
\text{DualCVRP}(t)\ \max\ &\sum_{i\in N_c}\sum_{v=\pi(i,t)}^{t-1} \lambda_{ivt}\,w_{iv}+K\,u_0\\
\text{s.t.}\ &\sum_{i\in N_c}\sum_{v=\pi(i,t)}^{t-1} a_{ivr}\,w_{iv}+u_0\le b_r,\quad \forall r\in\hat{\mathcal R}_t,\\
&w_{iv}\in\mathbb R,\quad \forall i\in N_c,\ v=\pi(i,t),\dots,t-1,\\
&u_0\le 0.
\end{aligned}
\label{eq:dual_cvrp_46}
\end{equation}

\begin{remark}
对任意 DualCVRP(t) 的可行解 $(w,u_0)$，其目标值都是 CVRP(t) 的\textbf{下界}（弱对偶）。
因此在 BMP 中添加 $\omega_t\ge \sum \lambda\,w+K u_0$ 形式的不等式一定是\textbf{安全的初始割}。
\end{remark}

\subsubsection*{B.2 初始割I：基于单周期松弛模型 $T1$ 的对偶（式(4-7)--(4-9)）}

为快速构造 DualCVRP(t) 的一个可行解，论文给出一个更容易（在实践中更“稳”）的松弛：
\begin{equation}
\begin{aligned}
(T1)\ \min\ &\sum_{r\in\hat{\mathcal R}_t} b_r\,\xi_r\\
\text{s.t.}\ &\sum_{r\in\hat{\mathcal R}_t} a_{ivr}\,\xi_r\ge 1,\quad \forall i\in N_c,\ v=\pi(i,t),\dots,t-1,\\
&\sum_{r\in\hat{\mathcal R}_t}\xi_r\le K,\\
&\xi_r\ge 0,\quad \forall r\in\hat{\mathcal R}_t.
\end{aligned}
\label{eq:t1_primal_47}
\end{equation}

其对偶（式(4-8)）为：
\begin{equation}
\begin{aligned}
(\text{Dual }T1)\ \max\ &\sum_{i\in N_c}\sum_{v=\pi(i,t)}^{t-1} w_{iv}+K u_0\\
\text{s.t.}\ &\sum_{i\in N_c}\sum_{v=\pi(i,t)}^{t-1} a_{ivr}\,w_{iv}+u_0\le b_r,\quad \forall r\in\hat{\mathcal R}_t,\\
&w_{iv}\ge 0,\quad \forall i\in N_c,\ v=\pi(i,t),\dots,t-1,\\
&u_0\le 0.
\end{aligned}
\label{eq:dual_t1_48}
\end{equation}

由于 \eqref{eq:dual_t1_48} 与 \eqref{eq:dual_cvrp_46} 的约束相同但更严格（$w_{iv}\ge 0$），所以 DualT1 的任意可行解也是 DualCVRP(t) 的可行解。
令 $(w^*,u_0^*)$ 为 DualT1 的任意一个可行解（通常取其最优解），则可加入初始割（式(4-9)）：
\begin{equation}
\omega_t\ \ge\ \sum_{i\in N_c}\sum_{v=\pi(i,t)}^{t-1} w_{iv}^*\,\lambda_{ivt}\ +\ K u_0^*.
\label{eq:initcut_t1_49}
\end{equation}

\paragraph{实现提示（怎样在代码里“拿到”$(w^*,u_0^*)$）}
\begin{itemize}[leftmargin=2em]
  \item 你可以直接解 \eqref{eq:dual_t1_48}（对偶）或 \eqref{eq:t1_primal_47}（原始）得到对偶变量。
  \item 由于 $\hat{\mathcal R}_t$ 很大，通常仍需要列生成：
    \begin{itemize}[leftmargin=2em]
      \item 给定 $(w,u_0)$，路径 $r$ 的对偶约束违反/定价等价于检查\;
      $\min_{r}\ \bigl(b_r-\sum_{i,v}a_{ivr}w_{iv}-u_0\bigr)$ 是否为负。
      \item 这与 CVRP 的“最短路/定价子问题”同源，可复用你已有的定价器。
    \end{itemize}
  \item 为了让系数更“有用”，论文建议采用 $\ge 1$（从而 $w_{iv}\ge 0$）而不是等式 $=1$（否则 $w_{iv}$ 无符号，容易给出较弱/不稳定的 cut 系数）。
\end{itemize}

\begin{remark}[$T1$ 初始割的全局有效性与强度权衡]
$T1$ 将覆盖约束从 $=\lambda_{ivt}$ 改为 $\ge 1$，使得 $T1$ 的求解\textbf{不依赖 $\lambda$ 的具体取值}——对所有情景节点 $(i,v)$ 均要求至少覆盖一次。这一特性带来两个直接后果：
\begin{itemize}
  \item \textbf{全局有效性}：割系数 $w^*_{iv}$ 对所有可能的 $\lambda$ 配置均成立，割 \eqref{eq:initcut_t1_49} 作为全局有效不等式可安全加入 BMP；
  \item \textbf{强度偏弱}：由于 $T1$ 忽略了访问序列约束（哪些弧实际被选中），其对偶解未能利用 $\lambda$ 结构，给出的下界通常比基于完整 $\lambda$ 信息的 $T2$ 割更松。
\end{itemize}
因此，$T1$ 割适合作为快速预热手段；若计算资源允许，建议进一步叠加 $T2$ 割以获得更紧的初始下界。
\end{remark}

\subsubsection*{B.3 初始割II：基于全局松弛模型 $T2$ 的对偶（式(4-10)--(4-12)）}

思想：$T1$ 忽略了很多来自“生产-库存-访问序列”的全局约束信息，得到的对偶解可能偏松；
因此构造一个更接近完整模型的 LP 松弛 $T2$，从其对偶可行解中抽取每期的 $(w^*_{ivt},u^*_{0t})$。

论文给出的 $T2$（式(4-10)）是“将第2.4节重构后的完整模型做 LP 松弛，并将路由部分用集合划分表示”：
\begin{equation}
\begin{aligned}
(T2)\ \min\ &\sum_{t\in T}\left(\sum_{r\in\hat{\mathcal R}_t} b_r\,\xi_{rt} + u p_t + f y_t + h_0 I_{0t}+h_p P_{0t}\right)
+\sum_{t\in T'}\sum_{i\in V_S}\sum_{v=\pi(i,t)}^{t-1} e_{ivt}\lambda_{ivt}\\
\text{s.t.}\ &\sum_{r\in\hat{\mathcal R}_t} a_{ivrt}\,\xi_{rt}=\lambda_{ivt},\quad \forall i\in N_c,\ t\in T,\ v=\pi(i,t),\dots,t-1,\\
&\sum_{r\in\hat{\mathcal R}_t}\xi_{rt}\le K,\quad \forall t\in T,\\
&0\le y_t\le 1,\ \forall t\in T,\qquad 0\le \lambda_{ivt}\le 1,\ \forall(i,v,t),\\
&\xi_{rt}\ge 0,\ \forall r\in\hat{\mathcal R}_t,\ t\in T,\\
&\text{并满足（重构模型中的其余线性约束），例如论文中的}\ (2\text{-}37)\sim(2\text{-}45),(2\text{-}60)\sim(2\text{-}62).
\end{aligned}
\label{eq:t2_lp_410}
\end{equation}

其对偶中，与 $\xi_{rt}$ 相关的对偶约束为（式(4-11)）：
\begin{equation}
\sum_{i\in N_c}\sum_{v=\pi(i,t)}^{t-1} a_{ivrt}\,w_{ivt}+u_{0t}\le b_r,\qquad \forall r\in\hat{\mathcal R}_t,\ t\in T.
\label{eq:t2_dual_constr_411}
\end{equation}

记 $(w^*,u_0^*)$ 为 $T2$ 的任意一个对偶可行解（通常取最优对偶解）。
论文给出的更紧初始割（式(4-12)）为：
\begin{equation}
\sum_{t\in T}\omega_t\ \ge\ \sum_{t\in T}\left(\sum_{i\in N_c}\sum_{v=\pi(i,t)}^{t-1} w_{ivt}^*\,\lambda_{ivt} + K u_{0t}^*\right).
\label{eq:initcut_t2_412}
\end{equation}

\paragraph{实现提示}
\begin{itemize}[leftmargin=2em]
  \item 与 $T1$ 相比，$T2$ 规模更大、求解更慢，但系数通常更紧。
  \item 论文形式写成“总和 cut” \eqref{eq:initcut_t2_412}（一次加一条），实现最简单：在开始迭代前解一次 $T2$ 的 LP，取对偶变量加一条不等式即可。
  \item 若你愿意，也可以把 \eqref{eq:initcut_t2_412} 拆成每期一条 $\omega_t\ge \sum w^*_{ivt}\lambda_{ivt}+K u^*_{0t}$（更紧），但这不属于论文原式，需自行做数值稳定性测试。
\end{itemize}

% ------------------------------------------------------------
% C. Correctness checks
% ------------------------------------------------------------
\subsection{C. 正确性检查（这些加速策略为什么“安全”）}

\paragraph{C.1 有效不等式（A.1--A.3）不会切掉最优整数解}
\begin{itemize}[leftmargin=2em]
  \item \eqref{eq:vi_start_collect_42}：
  若存在 $t'$ 使得 $P_{00}+I_{00}/k < \sum_{\tau=1}^{t'} d_\tau$，则仅靠初始库存无法满足到 $t'$ 的累积需求。
  在模型不允许欠交付（无 backlogging）且无外购原料的前提下，必须在 $t'$ 及之前产生至少一次原料“流入”，也就是至少一次收集访问（$z_{i\tau}=1$）。
  \item \eqref{eq:vi_totalcap_43}：
  任一可行路由方案在时期 $t$ 的总载重不超过 $KQ$，而左边正是主问题计划的总收集量，因此是必要条件。
  \item \eqref{eq:vi_bigitem_44}：
  若某条路径容量为 $Q$，则不可能同时装下两个需求都 $>Q/2$ 的客户，因此“大客户”的数量至多等于可用车辆数 $K$。
\end{itemize}

\paragraph{C.2 初始割（B.2--B.3）不会切掉最优解}
\begin{itemize}[leftmargin=2em]
  \item \eqref{eq:initcut_t1_49} 来自 DualCVRP(t) 的可行解：DualT1 的可行域是 DualCVRP(t) 的子集（仅额外要求 $w\ge0$），因此 $(w^*,u_0^*)$ 对 DualCVRP(t) 必可行。弱对偶给出 $\sum \lambda w^*+K u_0^*$ 是 CVRP(t) 的下界，从而 $\omega_t$ 作为 CVRP(t) 最优值必须不小于它。
  \item \eqref{eq:initcut_t2_412}：
  $T2$ 的对偶可行解在每个周期 $t$ 上的分量 $(w^*_{\cdot\cdot t},u^*_{0t})$ 满足与 DualCVRP(t) 同形的约束 \eqref{eq:t2_dual_constr_411}，故同样可作为每期 CVRP(t) 的对偶可行解并给出下界。将各期下界求和即可得到 \eqref{eq:initcut_t2_412}。
\end{itemize}

% ------------------------------------------------------------
% D. Implementation checklist
% ------------------------------------------------------------
\subsection{D. 代码实现 Checklist（建议按此顺序做）}

\begin{algorithm}[H]
\caption{把第4章加速策略接入 LBBD 的最小实现流程}
\DontPrintSemicolon
\KwIn{基础数据 + 第2.4节预计算参数 $\pi(i,t),g_{ivt}$（以及可选的 $e_{ivt}$）}
\KwIn{已实现的 LBBD 框架（BMP/BSP/列生成与割回调）}
\BlankLine
\textbf{(0) 主问题建模阶段：添加有效不等式}\;
计算 $t'$（若存在）并加入 \eqref{eq:vi_start_collect_42}\;
对每个 $t\in T$ 加入 \eqref{eq:vi_totalcap_43} 与 \eqref{eq:vi_bigitem_44}\;
\BlankLine
\textbf{(1) Benders 迭代前：生成初始割（可选，但推荐）}\;
\If{使用 $T1$ 初始割}{
对每个 $t\in T$ 解一次 DualT1（或解 $T1$ 并读出对偶变量），得到 $(w^*,u_0^*)$\;
向 BMP 加入 \eqref{eq:initcut_t1_49}\;
}
\If{使用 $T2$ 初始割}{
解一次全局 LP 松弛 $T2$ 并读出对偶可行解 $(w^*,u_0^*)$\;
向 BMP 加入 \eqref{eq:initcut_t2_412}\;
}
\BlankLine
\textbf{(2) 进入常规 LBBD 迭代：}\;
重复：解 BMP $\rightarrow$ 解各期 BSP(CVRP) $\rightarrow$ 生成并加入可行性/最优性割，直到收敛\;
\end{algorithm}

\end{document}
 或单独编译均可

\documentclass{article}
\usepackage{amsmath,amssymb,amsfonts,amsthm}
\usepackage{enumitem}
\usepackage{geometry}
\usepackage{mathrsfs}
\geometry{a4paper, margin=1in}
\usepackage{ctex} % 中文
\usepackage{hyperref}
\usepackage{float}
\usepackage[linesnumbered,ruled,vlined]{algorithm2e}

\newtheorem{remark}{注}
\newtheorem{definition}{定义}

\begin{document}

% ============================================================
%  Speed-up strategies for LBBD (Chapter 4)
%  1) Valid inequalities for BMP (master)
%  2) Warm-start initial cuts via dual-feasible solutions
% ============================================================

\section{LBBD 算法加速策略（第4章：有效不等式 + 初始割）}

\subsection{符号速查（只列本文件会用到的）}

\begin{itemize}[leftmargin=2em]
  \item 规划期：$T=\{1,\dots,l\}$。
  \item 收集点集合：$V_S$。在路由子问题中，客户集合写作 $N_c$（通常可取 $N_c=V_S$ 或“当期被访问的收集点”）。
  \item 车辆：可用车辆数上限 $K$，单车容量 $Q$。
  \item 主问题变量：
    \begin{itemize}[leftmargin=2em]
      \item $z_{it}\in\{0,1\}$：时期 $t$ 是否访问收集点 $i$。
      \item $\lambda_{ivt}\in\{0,1\}$：收集点 $i$ 的“上次访问时期为 $v$、本次访问时期为 $t$”的弧是否被选中（时间扩展网络/最短路式改写中的弧选择变量）。
      \item $\omega_t\ge 0$：时期 $t$ 的最优路径成本（由子问题给出，用 Benders 割约束它）。
    \end{itemize}
  \item 预计算参数（来自第2.4节改写模型；加速策略只\textbf{引用}它们，不重复推导）：
    \begin{itemize}[leftmargin=2em]
      \item $\pi(i,t)$：在“本次访问 $t$”条件下，\emph{可行的}最早上次访问时期下界（由收集点容量等逻辑得到）。
      \item $g_{ivt}$：若选择弧 $(v,t)$，即 $i$ 上次在 $v$ 被访问、当前在 $t$ 被访问，则时期 $t$ 从 $i$ 需要收走的\textbf{全量}（累计）收集量。
      \item （若实现 $T2$）$e_{ivt}$：与弧 $(v,t)$ 关联的收集点库存持有成本（或其它在第2.4节中被“弧化”的成本项）。
    \end{itemize}
\end{itemize}

\begin{remark}
本文件只覆盖第4章“加速方法”，默认你已经按 \texttt{lbbd\_all.tex} 中的说明实现了：
(1) BMP（主问题）中用 $\omega_t$ 代替显式路由变量；(2) BSP（子问题）用集合划分/列生成求解 CVRP(t) 并生成 Benders 割。
\end{remark}

% ------------------------------------------------------------
% A. Valid inequalities (add to BMP)
% ------------------------------------------------------------
\subsection{A. 主问题（BMP）有效不等式：直接加到主问题即可}

\subsubsection*{A.1 最小收集启动时间约束（式(4-1)--(4-2)）}

定义一个最早必须启动收集或生产的时期 $t'$：
\begin{equation}
t'=\arg\min_{1\le t\le l}\left\{\,P_{00}+\frac{I_{00}}{k}<\sum_{\tau=1}^{t} d_{\tau}\right\}.
\label{eq:tprime_def_41}
\end{equation}

若这样的 $t'$ 存在，则到 $t'$ 期末必须至少发生一次收集活动（否则没有新增原料来源，无法满足累积需求）：
\begin{equation}
\sum_{\tau=1}^{t'}\ \sum_{i\in V_S} z_{i\tau}\ \ge\ 1.
\label{eq:vi_start_collect_42}
\end{equation}

\paragraph{实现提示}
\begin{itemize}[leftmargin=2em]
  \item 这是一个\textbf{一次性}约束：只需在建模时计算 $t'$ 并加入一次 \eqref{eq:vi_start_collect_42}。
  \item 若不存在 $t'$（即初始库存 $P_{00}+I_{00}/k$ 足以覆盖全周期累积需求），则不要添加该约束。
\end{itemize}

\subsubsection*{A.2 周期车辆总容量约束（式(4-3)）}

时期 $t$ 的总计划收集量可由 $\lambda$ 的选择唯一确定（全量收集）：
\[
\text{Pickup}_t \;=\; \sum_{i\in V_S}\ \sum_{v=\pi(i,t)}^{t-1} g_{ivt}\,\lambda_{ivt}.
\]
由于每期最多使用 $K$ 辆容量为 $Q$ 的车，必要条件为：
\begin{equation}
\sum_{i\in V_S}\ \sum_{v=\pi(i,t)}^{t-1} g_{ivt}\,\lambda_{ivt}\ \le\ K\cdot Q,\qquad \forall t\in T.
\label{eq:vi_totalcap_43}
\end{equation}

\paragraph{实现提示}
\begin{itemize}[leftmargin=2em]
  \item 这是\textbf{按期}约束：对所有 $t\in T$ 添加 \eqref{eq:vi_totalcap_43}。
  \item 该不等式是可行性的\textbf{必要条件}（不是充分条件），但能显著收紧主问题的 LP 松弛。
\end{itemize}

\subsubsection*{A.3 大收集量客户约束（式(4-4)）}

若某个收集点在时期 $t$ 的收集量超过半车容量（$g_{ivt}>Q/2$），则同一条路径不可能再服务另一个同样“大”的收集点。
因此，在时期 $t$，满足 $g_{ivt}>Q/2$ 的被选中弧数量不能超过 $K$：
\begin{equation}
\sum_{i\in V_S}\ \sum_{\substack{v=\pi(i,t)\\ g_{ivt}>Q/2}}^{t-1}\ \lambda_{ivt}\ \le\ K,\qquad \forall t\in T.
\label{eq:vi_bigitem_44}
\end{equation}

\paragraph{实现提示}
\begin{itemize}[leftmargin=2em]
  \item 这同样是\textbf{按期}约束：对每个 $t$ 建一个线性约束，内部只枚举满足 $g_{ivt}>Q/2$ 的 $(i,v)$。
  \item 若存在 $g_{ivt}>Q$，则该弧本身会导致 CVRP(t) 不可行，建议在预处理阶段直接\textbf{删除}该弧或在主问题中禁用它（比加 \eqref{eq:vi_bigitem_44} 更“直接”）。
\end{itemize}

% ------------------------------------------------------------
% B. Initial cuts (warm start)
% ------------------------------------------------------------
\subsection{B. 主问题（BMP）初始割：Benders 迭代前的“预热”}

动机：LBBD 初期，$\omega_t$ 在主问题里可能被压到非常小（甚至 0），导致第一次/前几次迭代的主问题解偏离真实可行解，从而迭代次数增多。初始割为 $\omega_t$ 提供一个\textbf{有效下界}，加速收敛。

\subsubsection*{B.1 访问决策未知的 CVRP(t) 集合划分模型（式(4-5)）}

设 $\hat{\mathcal R}_t$ 为时期 $t$ 的所有可行路径集合（通过列生成隐式表示），路径 $r$ 的成本为 $b_r$，其对“客户-需求情景”$(i,v)$ 的覆盖系数为 $a_{ivr}\in\{0,1\}$。
则 CVRP(t) 可写为：
\begin{equation}
\begin{aligned}
\text{CVRP}(t)\;\; \omega_t=\min\ &\sum_{r\in\hat{\mathcal R}_t} b_r\,\xi_r\\
\text{s.t.}\ &\sum_{r\in\hat{\mathcal R}_t} a_{ivr}\,\xi_r=\lambda_{ivt},\quad \forall i\in N_c,\ v=\pi(i,t),\dots,t-1,\\
&\sum_{r\in\hat{\mathcal R}_t}\xi_r\le K,\\
&\xi_r\in\{0,1\},\quad \forall r\in\hat{\mathcal R}_t.
\end{aligned}
\label{eq:cvrp_sp_45}
\end{equation}

其 LP 松弛对偶为（式(4-6)）：
\begin{equation}
\begin{aligned}
\text{DualCVRP}(t)\ \max\ &\sum_{i\in N_c}\sum_{v=\pi(i,t)}^{t-1} \lambda_{ivt}\,w_{iv}+K\,u_0\\
\text{s.t.}\ &\sum_{i\in N_c}\sum_{v=\pi(i,t)}^{t-1} a_{ivr}\,w_{iv}+u_0\le b_r,\quad \forall r\in\hat{\mathcal R}_t,\\
&w_{iv}\in\mathbb R,\quad \forall i\in N_c,\ v=\pi(i,t),\dots,t-1,\\
&u_0\le 0.
\end{aligned}
\label{eq:dual_cvrp_46}
\end{equation}

\begin{remark}
对任意 DualCVRP(t) 的可行解 $(w,u_0)$，其目标值都是 CVRP(t) 的\textbf{下界}（弱对偶）。
因此在 BMP 中添加 $\omega_t\ge \sum \lambda\,w+K u_0$ 形式的不等式一定是\textbf{安全的初始割}。
\end{remark}

\subsubsection*{B.2 初始割I：基于单周期松弛模型 $T1$ 的对偶（式(4-7)--(4-9)）}

为快速构造 DualCVRP(t) 的一个可行解，论文给出一个更容易（在实践中更“稳”）的松弛：
\begin{equation}
\begin{aligned}
(T1)\ \min\ &\sum_{r\in\hat{\mathcal R}_t} b_r\,\xi_r\\
\text{s.t.}\ &\sum_{r\in\hat{\mathcal R}_t} a_{ivr}\,\xi_r\ge 1,\quad \forall i\in N_c,\ v=\pi(i,t),\dots,t-1,\\
&\sum_{r\in\hat{\mathcal R}_t}\xi_r\le K,\\
&\xi_r\ge 0,\quad \forall r\in\hat{\mathcal R}_t.
\end{aligned}
\label{eq:t1_primal_47}
\end{equation}

其对偶（式(4-8)）为：
\begin{equation}
\begin{aligned}
(\text{Dual }T1)\ \max\ &\sum_{i\in N_c}\sum_{v=\pi(i,t)}^{t-1} w_{iv}+K u_0\\
\text{s.t.}\ &\sum_{i\in N_c}\sum_{v=\pi(i,t)}^{t-1} a_{ivr}\,w_{iv}+u_0\le b_r,\quad \forall r\in\hat{\mathcal R}_t,\\
&w_{iv}\ge 0,\quad \forall i\in N_c,\ v=\pi(i,t),\dots,t-1,\\
&u_0\le 0.
\end{aligned}
\label{eq:dual_t1_48}
\end{equation}

由于 \eqref{eq:dual_t1_48} 与 \eqref{eq:dual_cvrp_46} 的约束相同但更严格（$w_{iv}\ge 0$），所以 DualT1 的任意可行解也是 DualCVRP(t) 的可行解。
令 $(w^*,u_0^*)$ 为 DualT1 的任意一个可行解（通常取其最优解），则可加入初始割（式(4-9)）：
\begin{equation}
\omega_t\ \ge\ \sum_{i\in N_c}\sum_{v=\pi(i,t)}^{t-1} w_{iv}^*\,\lambda_{ivt}\ +\ K u_0^*.
\label{eq:initcut_t1_49}
\end{equation}

\paragraph{实现提示（怎样在代码里“拿到”$(w^*,u_0^*)$）}
\begin{itemize}[leftmargin=2em]
  \item 你可以直接解 \eqref{eq:dual_t1_48}（对偶）或 \eqref{eq:t1_primal_47}（原始）得到对偶变量。
  \item 由于 $\hat{\mathcal R}_t$ 很大，通常仍需要列生成：
    \begin{itemize}[leftmargin=2em]
      \item 给定 $(w,u_0)$，路径 $r$ 的对偶约束违反/定价等价于检查\;
      $\min_{r}\ \bigl(b_r-\sum_{i,v}a_{ivr}w_{iv}-u_0\bigr)$ 是否为负。
      \item 这与 CVRP 的“最短路/定价子问题”同源，可复用你已有的定价器。
    \end{itemize}
  \item 为了让系数更“有用”，论文建议采用 $\ge 1$（从而 $w_{iv}\ge 0$）而不是等式 $=1$（否则 $w_{iv}$ 无符号，容易给出较弱/不稳定的 cut 系数）。
\end{itemize}

\begin{remark}[$T1$ 初始割的全局有效性与强度权衡]
$T1$ 将覆盖约束从 $=\lambda_{ivt}$ 改为 $\ge 1$，使得 $T1$ 的求解\textbf{不依赖 $\lambda$ 的具体取值}——对所有情景节点 $(i,v)$ 均要求至少覆盖一次。这一特性带来两个直接后果：
\begin{itemize}
  \item \textbf{全局有效性}：割系数 $w^*_{iv}$ 对所有可能的 $\lambda$ 配置均成立，割 \eqref{eq:initcut_t1_49} 作为全局有效不等式可安全加入 BMP；
  \item \textbf{强度偏弱}：由于 $T1$ 忽略了访问序列约束（哪些弧实际被选中），其对偶解未能利用 $\lambda$ 结构，给出的下界通常比基于完整 $\lambda$ 信息的 $T2$ 割更松。
\end{itemize}
因此，$T1$ 割适合作为快速预热手段；若计算资源允许，建议进一步叠加 $T2$ 割以获得更紧的初始下界。
\end{remark}

\subsubsection*{B.3 初始割II：基于全局松弛模型 $T2$ 的对偶（式(4-10)--(4-12)）}

思想：$T1$ 忽略了很多来自“生产-库存-访问序列”的全局约束信息，得到的对偶解可能偏松；
因此构造一个更接近完整模型的 LP 松弛 $T2$，从其对偶可行解中抽取每期的 $(w^*_{ivt},u^*_{0t})$。

论文给出的 $T2$（式(4-10)）是“将第2.4节重构后的完整模型做 LP 松弛，并将路由部分用集合划分表示”：
\begin{equation}
\begin{aligned}
(T2)\ \min\ &\sum_{t\in T}\left(\sum_{r\in\hat{\mathcal R}_t} b_r\,\xi_{rt} + u p_t + f y_t + h_0 I_{0t}+h_p P_{0t}\right)
+\sum_{t\in T'}\sum_{i\in V_S}\sum_{v=\pi(i,t)}^{t-1} e_{ivt}\lambda_{ivt}\\
\text{s.t.}\ &\sum_{r\in\hat{\mathcal R}_t} a_{ivrt}\,\xi_{rt}=\lambda_{ivt},\quad \forall i\in N_c,\ t\in T,\ v=\pi(i,t),\dots,t-1,\\
&\sum_{r\in\hat{\mathcal R}_t}\xi_{rt}\le K,\quad \forall t\in T,\\
&0\le y_t\le 1,\ \forall t\in T,\qquad 0\le \lambda_{ivt}\le 1,\ \forall(i,v,t),\\
&\xi_{rt}\ge 0,\ \forall r\in\hat{\mathcal R}_t,\ t\in T,\\
&\text{并满足（重构模型中的其余线性约束），例如论文中的}\ (2\text{-}37)\sim(2\text{-}45),(2\text{-}60)\sim(2\text{-}62).
\end{aligned}
\label{eq:t2_lp_410}
\end{equation}

其对偶中，与 $\xi_{rt}$ 相关的对偶约束为（式(4-11)）：
\begin{equation}
\sum_{i\in N_c}\sum_{v=\pi(i,t)}^{t-1} a_{ivrt}\,w_{ivt}+u_{0t}\le b_r,\qquad \forall r\in\hat{\mathcal R}_t,\ t\in T.
\label{eq:t2_dual_constr_411}
\end{equation}

记 $(w^*,u_0^*)$ 为 $T2$ 的任意一个对偶可行解（通常取最优对偶解）。
论文给出的更紧初始割（式(4-12)）为：
\begin{equation}
\sum_{t\in T}\omega_t\ \ge\ \sum_{t\in T}\left(\sum_{i\in N_c}\sum_{v=\pi(i,t)}^{t-1} w_{ivt}^*\,\lambda_{ivt} + K u_{0t}^*\right).
\label{eq:initcut_t2_412}
\end{equation}

\paragraph{实现提示}
\begin{itemize}[leftmargin=2em]
  \item 与 $T1$ 相比，$T2$ 规模更大、求解更慢，但系数通常更紧。
  \item 论文形式写成“总和 cut” \eqref{eq:initcut_t2_412}（一次加一条），实现最简单：在开始迭代前解一次 $T2$ 的 LP，取对偶变量加一条不等式即可。
  \item 若你愿意，也可以把 \eqref{eq:initcut_t2_412} 拆成每期一条 $\omega_t\ge \sum w^*_{ivt}\lambda_{ivt}+K u^*_{0t}$（更紧），但这不属于论文原式，需自行做数值稳定性测试。
\end{itemize}

% ------------------------------------------------------------
% C. Correctness checks
% ------------------------------------------------------------
\subsection{C. 正确性检查（这些加速策略为什么“安全”）}

\paragraph{C.1 有效不等式（A.1--A.3）不会切掉最优整数解}
\begin{itemize}[leftmargin=2em]
  \item \eqref{eq:vi_start_collect_42}：
  若存在 $t'$ 使得 $P_{00}+I_{00}/k < \sum_{\tau=1}^{t'} d_\tau$，则仅靠初始库存无法满足到 $t'$ 的累积需求。
  在模型不允许欠交付（无 backlogging）且无外购原料的前提下，必须在 $t'$ 及之前产生至少一次原料“流入”，也就是至少一次收集访问（$z_{i\tau}=1$）。
  \item \eqref{eq:vi_totalcap_43}：
  任一可行路由方案在时期 $t$ 的总载重不超过 $KQ$，而左边正是主问题计划的总收集量，因此是必要条件。
  \item \eqref{eq:vi_bigitem_44}：
  若某条路径容量为 $Q$，则不可能同时装下两个需求都 $>Q/2$ 的客户，因此“大客户”的数量至多等于可用车辆数 $K$。
\end{itemize}

\paragraph{C.2 初始割（B.2--B.3）不会切掉最优解}
\begin{itemize}[leftmargin=2em]
  \item \eqref{eq:initcut_t1_49} 来自 DualCVRP(t) 的可行解：DualT1 的可行域是 DualCVRP(t) 的子集（仅额外要求 $w\ge0$），因此 $(w^*,u_0^*)$ 对 DualCVRP(t) 必可行。弱对偶给出 $\sum \lambda w^*+K u_0^*$ 是 CVRP(t) 的下界，从而 $\omega_t$ 作为 CVRP(t) 最优值必须不小于它。
  \item \eqref{eq:initcut_t2_412}：
  $T2$ 的对偶可行解在每个周期 $t$ 上的分量 $(w^*_{\cdot\cdot t},u^*_{0t})$ 满足与 DualCVRP(t) 同形的约束 \eqref{eq:t2_dual_constr_411}，故同样可作为每期 CVRP(t) 的对偶可行解并给出下界。将各期下界求和即可得到 \eqref{eq:initcut_t2_412}。
\end{itemize}

% ------------------------------------------------------------
% D. Implementation checklist
% ------------------------------------------------------------
\subsection{D. 代码实现 Checklist（建议按此顺序做）}

\begin{algorithm}[H]
\caption{把第4章加速策略接入 LBBD 的最小实现流程}
\DontPrintSemicolon
\KwIn{基础数据 + 第2.4节预计算参数 $\pi(i,t),g_{ivt}$（以及可选的 $e_{ivt}$）}
\KwIn{已实现的 LBBD 框架（BMP/BSP/列生成与割回调）}
\BlankLine
\textbf{(0) 主问题建模阶段：添加有效不等式}\;
计算 $t'$（若存在）并加入 \eqref{eq:vi_start_collect_42}\;
对每个 $t\in T$ 加入 \eqref{eq:vi_totalcap_43} 与 \eqref{eq:vi_bigitem_44}\;
\BlankLine
\textbf{(1) Benders 迭代前：生成初始割（可选，但推荐）}\;
\If{使用 $T1$ 初始割}{
对每个 $t\in T$ 解一次 DualT1（或解 $T1$ 并读出对偶变量），得到 $(w^*,u_0^*)$\;
向 BMP 加入 \eqref{eq:initcut_t1_49}\;
}
\If{使用 $T2$ 初始割}{
解一次全局 LP 松弛 $T2$ 并读出对偶可行解 $(w^*,u_0^*)$\;
向 BMP 加入 \eqref{eq:initcut_t2_412}\;
}
\BlankLine
\textbf{(2) 进入常规 LBBD 迭代：}\;
重复：解 BMP $\rightarrow$ 解各期 BSP(CVRP) $\rightarrow$ 生成并加入可行性/最优性割，直到收敛\;
\end{algorithm}

\end{document}
 或单独编译均可

\documentclass{article}
\usepackage{amsmath,amssymb,amsfonts,amsthm}
\usepackage{enumitem}
\usepackage{geometry}
\usepackage{mathrsfs}
\geometry{a4paper, margin=1in}
\usepackage{ctex} % 中文
\usepackage{hyperref}
\usepackage{float}
\usepackage[linesnumbered,ruled,vlined]{algorithm2e}

\newtheorem{remark}{注}
\newtheorem{definition}{定义}

\begin{document}

% ============================================================
%  Speed-up strategies for LBBD (Chapter 4)
%  1) Valid inequalities for BMP (master)
%  2) Warm-start initial cuts via dual-feasible solutions
% ============================================================

\section{LBBD 算法加速策略（第4章：有效不等式 + 初始割）}

\subsection{符号速查（只列本文件会用到的）}

\begin{itemize}[leftmargin=2em]
  \item 规划期：$T=\{1,\dots,l\}$。
  \item 收集点集合：$V_S$。在路由子问题中，客户集合写作 $N_c$（通常可取 $N_c=V_S$ 或“当期被访问的收集点”）。
  \item 车辆：可用车辆数上限 $K$，单车容量 $Q$。
  \item 主问题变量：
    \begin{itemize}[leftmargin=2em]
      \item $z_{it}\in\{0,1\}$：时期 $t$ 是否访问收集点 $i$。
      \item $\lambda_{ivt}\in\{0,1\}$：收集点 $i$ 的“上次访问时期为 $v$、本次访问时期为 $t$”的弧是否被选中（时间扩展网络/最短路式改写中的弧选择变量）。
      \item $\omega_t\ge 0$：时期 $t$ 的最优路径成本（由子问题给出，用 Benders 割约束它）。
    \end{itemize}
  \item 预计算参数（来自第2.4节改写模型；加速策略只\textbf{引用}它们，不重复推导）：
    \begin{itemize}[leftmargin=2em]
      \item $\pi(i,t)$：在“本次访问 $t$”条件下，\emph{可行的}最早上次访问时期下界（由收集点容量等逻辑得到）。
      \item $g_{ivt}$：若选择弧 $(v,t)$，即 $i$ 上次在 $v$ 被访问、当前在 $t$ 被访问，则时期 $t$ 从 $i$ 需要收走的\textbf{全量}（累计）收集量。
      \item （若实现 $T2$）$e_{ivt}$：与弧 $(v,t)$ 关联的收集点库存持有成本（或其它在第2.4节中被“弧化”的成本项）。
    \end{itemize}
\end{itemize}

\begin{remark}
本文件只覆盖第4章“加速方法”，默认你已经按 \texttt{lbbd\_all.tex} 中的说明实现了：
(1) BMP（主问题）中用 $\omega_t$ 代替显式路由变量；(2) BSP（子问题）用集合划分/列生成求解 CVRP(t) 并生成 Benders 割。
\end{remark}

% ------------------------------------------------------------
% A. Valid inequalities (add to BMP)
% ------------------------------------------------------------
\subsection{A. 主问题（BMP）有效不等式：直接加到主问题即可}

\subsubsection*{A.1 最小收集启动时间约束（式(4-1)--(4-2)）}

定义一个最早必须启动收集或生产的时期 $t'$：
\begin{equation}
t'=\arg\min_{1\le t\le l}\left\{\,P_{00}+\frac{I_{00}}{k}<\sum_{\tau=1}^{t} d_{\tau}\right\}.
\label{eq:tprime_def_41}
\end{equation}

若这样的 $t'$ 存在，则到 $t'$ 期末必须至少发生一次收集活动（否则没有新增原料来源，无法满足累积需求）：
\begin{equation}
\sum_{\tau=1}^{t'}\ \sum_{i\in V_S} z_{i\tau}\ \ge\ 1.
\label{eq:vi_start_collect_42}
\end{equation}

\paragraph{实现提示}
\begin{itemize}[leftmargin=2em]
  \item 这是一个\textbf{一次性}约束：只需在建模时计算 $t'$ 并加入一次 \eqref{eq:vi_start_collect_42}。
  \item 若不存在 $t'$（即初始库存 $P_{00}+I_{00}/k$ 足以覆盖全周期累积需求），则不要添加该约束。
\end{itemize}

\subsubsection*{A.2 周期车辆总容量约束（式(4-3)）}

时期 $t$ 的总计划收集量可由 $\lambda$ 的选择唯一确定（全量收集）：
\[
\text{Pickup}_t \;=\; \sum_{i\in V_S}\ \sum_{v=\pi(i,t)}^{t-1} g_{ivt}\,\lambda_{ivt}.
\]
由于每期最多使用 $K$ 辆容量为 $Q$ 的车，必要条件为：
\begin{equation}
\sum_{i\in V_S}\ \sum_{v=\pi(i,t)}^{t-1} g_{ivt}\,\lambda_{ivt}\ \le\ K\cdot Q,\qquad \forall t\in T.
\label{eq:vi_totalcap_43}
\end{equation}

\paragraph{实现提示}
\begin{itemize}[leftmargin=2em]
  \item 这是\textbf{按期}约束：对所有 $t\in T$ 添加 \eqref{eq:vi_totalcap_43}。
  \item 该不等式是可行性的\textbf{必要条件}（不是充分条件），但能显著收紧主问题的 LP 松弛。
\end{itemize}

\subsubsection*{A.3 大收集量客户约束（式(4-4)）}

若某个收集点在时期 $t$ 的收集量超过半车容量（$g_{ivt}>Q/2$），则同一条路径不可能再服务另一个同样“大”的收集点。
因此，在时期 $t$，满足 $g_{ivt}>Q/2$ 的被选中弧数量不能超过 $K$：
\begin{equation}
\sum_{i\in V_S}\ \sum_{\substack{v=\pi(i,t)\\ g_{ivt}>Q/2}}^{t-1}\ \lambda_{ivt}\ \le\ K,\qquad \forall t\in T.
\label{eq:vi_bigitem_44}
\end{equation}

\paragraph{实现提示}
\begin{itemize}[leftmargin=2em]
  \item 这同样是\textbf{按期}约束：对每个 $t$ 建一个线性约束，内部只枚举满足 $g_{ivt}>Q/2$ 的 $(i,v)$。
  \item 若存在 $g_{ivt}>Q$，则该弧本身会导致 CVRP(t) 不可行，建议在预处理阶段直接\textbf{删除}该弧或在主问题中禁用它（比加 \eqref{eq:vi_bigitem_44} 更“直接”）。
\end{itemize}

% ------------------------------------------------------------
% B. Initial cuts (warm start)
% ------------------------------------------------------------
\subsection{B. 主问题（BMP）初始割：Benders 迭代前的“预热”}

动机：LBBD 初期，$\omega_t$ 在主问题里可能被压到非常小（甚至 0），导致第一次/前几次迭代的主问题解偏离真实可行解，从而迭代次数增多。初始割为 $\omega_t$ 提供一个\textbf{有效下界}，加速收敛。

\subsubsection*{B.1 访问决策未知的 CVRP(t) 集合划分模型（式(4-5)）}

设 $\hat{\mathcal R}_t$ 为时期 $t$ 的所有可行路径集合（通过列生成隐式表示），路径 $r$ 的成本为 $b_r$，其对“客户-需求情景”$(i,v)$ 的覆盖系数为 $a_{ivr}\in\{0,1\}$。
则 CVRP(t) 可写为：
\begin{equation}
\begin{aligned}
\text{CVRP}(t)\;\; \omega_t=\min\ &\sum_{r\in\hat{\mathcal R}_t} b_r\,\xi_r\\
\text{s.t.}\ &\sum_{r\in\hat{\mathcal R}_t} a_{ivr}\,\xi_r=\lambda_{ivt},\quad \forall i\in N_c,\ v=\pi(i,t),\dots,t-1,\\
&\sum_{r\in\hat{\mathcal R}_t}\xi_r\le K,\\
&\xi_r\in\{0,1\},\quad \forall r\in\hat{\mathcal R}_t.
\end{aligned}
\label{eq:cvrp_sp_45}
\end{equation}

其 LP 松弛对偶为（式(4-6)）：
\begin{equation}
\begin{aligned}
\text{DualCVRP}(t)\ \max\ &\sum_{i\in N_c}\sum_{v=\pi(i,t)}^{t-1} \lambda_{ivt}\,w_{iv}+K\,u_0\\
\text{s.t.}\ &\sum_{i\in N_c}\sum_{v=\pi(i,t)}^{t-1} a_{ivr}\,w_{iv}+u_0\le b_r,\quad \forall r\in\hat{\mathcal R}_t,\\
&w_{iv}\in\mathbb R,\quad \forall i\in N_c,\ v=\pi(i,t),\dots,t-1,\\
&u_0\le 0.
\end{aligned}
\label{eq:dual_cvrp_46}
\end{equation}

\begin{remark}
对任意 DualCVRP(t) 的可行解 $(w,u_0)$，其目标值都是 CVRP(t) 的\textbf{下界}（弱对偶）。
因此在 BMP 中添加 $\omega_t\ge \sum \lambda\,w+K u_0$ 形式的不等式一定是\textbf{安全的初始割}。
\end{remark}

\subsubsection*{B.2 初始割I：基于单周期松弛模型 $T1$ 的对偶（式(4-7)--(4-9)）}

为快速构造 DualCVRP(t) 的一个可行解，论文给出一个更容易（在实践中更“稳”）的松弛：
\begin{equation}
\begin{aligned}
(T1)\ \min\ &\sum_{r\in\hat{\mathcal R}_t} b_r\,\xi_r\\
\text{s.t.}\ &\sum_{r\in\hat{\mathcal R}_t} a_{ivr}\,\xi_r\ge 1,\quad \forall i\in N_c,\ v=\pi(i,t),\dots,t-1,\\
&\sum_{r\in\hat{\mathcal R}_t}\xi_r\le K,\\
&\xi_r\ge 0,\quad \forall r\in\hat{\mathcal R}_t.
\end{aligned}
\label{eq:t1_primal_47}
\end{equation}

其对偶（式(4-8)）为：
\begin{equation}
\begin{aligned}
(\text{Dual }T1)\ \max\ &\sum_{i\in N_c}\sum_{v=\pi(i,t)}^{t-1} w_{iv}+K u_0\\
\text{s.t.}\ &\sum_{i\in N_c}\sum_{v=\pi(i,t)}^{t-1} a_{ivr}\,w_{iv}+u_0\le b_r,\quad \forall r\in\hat{\mathcal R}_t,\\
&w_{iv}\ge 0,\quad \forall i\in N_c,\ v=\pi(i,t),\dots,t-1,\\
&u_0\le 0.
\end{aligned}
\label{eq:dual_t1_48}
\end{equation}

由于 \eqref{eq:dual_t1_48} 与 \eqref{eq:dual_cvrp_46} 的约束相同但更严格（$w_{iv}\ge 0$），所以 DualT1 的任意可行解也是 DualCVRP(t) 的可行解。
令 $(w^*,u_0^*)$ 为 DualT1 的任意一个可行解（通常取其最优解），则可加入初始割（式(4-9)）：
\begin{equation}
\omega_t\ \ge\ \sum_{i\in N_c}\sum_{v=\pi(i,t)}^{t-1} w_{iv}^*\,\lambda_{ivt}\ +\ K u_0^*.
\label{eq:initcut_t1_49}
\end{equation}

\paragraph{实现提示（怎样在代码里“拿到”$(w^*,u_0^*)$）}
\begin{itemize}[leftmargin=2em]
  \item 你可以直接解 \eqref{eq:dual_t1_48}（对偶）或 \eqref{eq:t1_primal_47}（原始）得到对偶变量。
  \item 由于 $\hat{\mathcal R}_t$ 很大，通常仍需要列生成：
    \begin{itemize}[leftmargin=2em]
      \item 给定 $(w,u_0)$，路径 $r$ 的对偶约束违反/定价等价于检查\;
      $\min_{r}\ \bigl(b_r-\sum_{i,v}a_{ivr}w_{iv}-u_0\bigr)$ 是否为负。
      \item 这与 CVRP 的“最短路/定价子问题”同源，可复用你已有的定价器。
    \end{itemize}
  \item 为了让系数更“有用”，论文建议采用 $\ge 1$（从而 $w_{iv}\ge 0$）而不是等式 $=1$（否则 $w_{iv}$ 无符号，容易给出较弱/不稳定的 cut 系数）。
\end{itemize}

\begin{remark}[$T1$ 初始割的全局有效性与强度权衡]
$T1$ 将覆盖约束从 $=\lambda_{ivt}$ 改为 $\ge 1$，使得 $T1$ 的求解\textbf{不依赖 $\lambda$ 的具体取值}——对所有情景节点 $(i,v)$ 均要求至少覆盖一次。这一特性带来两个直接后果：
\begin{itemize}
  \item \textbf{全局有效性}：割系数 $w^*_{iv}$ 对所有可能的 $\lambda$ 配置均成立，割 \eqref{eq:initcut_t1_49} 作为全局有效不等式可安全加入 BMP；
  \item \textbf{强度偏弱}：由于 $T1$ 忽略了访问序列约束（哪些弧实际被选中），其对偶解未能利用 $\lambda$ 结构，给出的下界通常比基于完整 $\lambda$ 信息的 $T2$ 割更松。
\end{itemize}
因此，$T1$ 割适合作为快速预热手段；若计算资源允许，建议进一步叠加 $T2$ 割以获得更紧的初始下界。
\end{remark}

\subsubsection*{B.3 初始割II：基于全局松弛模型 $T2$ 的对偶（式(4-10)--(4-12)）}

思想：$T1$ 忽略了很多来自“生产-库存-访问序列”的全局约束信息，得到的对偶解可能偏松；
因此构造一个更接近完整模型的 LP 松弛 $T2$，从其对偶可行解中抽取每期的 $(w^*_{ivt},u^*_{0t})$。

论文给出的 $T2$（式(4-10)）是“将第2.4节重构后的完整模型做 LP 松弛，并将路由部分用集合划分表示”：
\begin{equation}
\begin{aligned}
(T2)\ \min\ &\sum_{t\in T}\left(\sum_{r\in\hat{\mathcal R}_t} b_r\,\xi_{rt} + u p_t + f y_t + h_0 I_{0t}+h_p P_{0t}\right)
+\sum_{t\in T'}\sum_{i\in V_S}\sum_{v=\pi(i,t)}^{t-1} e_{ivt}\lambda_{ivt}\\
\text{s.t.}\ &\sum_{r\in\hat{\mathcal R}_t} a_{ivrt}\,\xi_{rt}=\lambda_{ivt},\quad \forall i\in N_c,\ t\in T,\ v=\pi(i,t),\dots,t-1,\\
&\sum_{r\in\hat{\mathcal R}_t}\xi_{rt}\le K,\quad \forall t\in T,\\
&0\le y_t\le 1,\ \forall t\in T,\qquad 0\le \lambda_{ivt}\le 1,\ \forall(i,v,t),\\
&\xi_{rt}\ge 0,\ \forall r\in\hat{\mathcal R}_t,\ t\in T,\\
&\text{并满足（重构模型中的其余线性约束），例如论文中的}\ (2\text{-}37)\sim(2\text{-}45),(2\text{-}60)\sim(2\text{-}62).
\end{aligned}
\label{eq:t2_lp_410}
\end{equation}

其对偶中，与 $\xi_{rt}$ 相关的对偶约束为（式(4-11)）：
\begin{equation}
\sum_{i\in N_c}\sum_{v=\pi(i,t)}^{t-1} a_{ivrt}\,w_{ivt}+u_{0t}\le b_r,\qquad \forall r\in\hat{\mathcal R}_t,\ t\in T.
\label{eq:t2_dual_constr_411}
\end{equation}

记 $(w^*,u_0^*)$ 为 $T2$ 的任意一个对偶可行解（通常取最优对偶解）。
论文给出的更紧初始割（式(4-12)）为：
\begin{equation}
\sum_{t\in T}\omega_t\ \ge\ \sum_{t\in T}\left(\sum_{i\in N_c}\sum_{v=\pi(i,t)}^{t-1} w_{ivt}^*\,\lambda_{ivt} + K u_{0t}^*\right).
\label{eq:initcut_t2_412}
\end{equation}

\paragraph{实现提示}
\begin{itemize}[leftmargin=2em]
  \item 与 $T1$ 相比，$T2$ 规模更大、求解更慢，但系数通常更紧。
  \item 论文形式写成“总和 cut” \eqref{eq:initcut_t2_412}（一次加一条），实现最简单：在开始迭代前解一次 $T2$ 的 LP，取对偶变量加一条不等式即可。
  \item 若你愿意，也可以把 \eqref{eq:initcut_t2_412} 拆成每期一条 $\omega_t\ge \sum w^*_{ivt}\lambda_{ivt}+K u^*_{0t}$（更紧），但这不属于论文原式，需自行做数值稳定性测试。
\end{itemize}

% ------------------------------------------------------------
% C. Correctness checks
% ------------------------------------------------------------
\subsection{C. 正确性检查（这些加速策略为什么“安全”）}

\paragraph{C.1 有效不等式（A.1--A.3）不会切掉最优整数解}
\begin{itemize}[leftmargin=2em]
  \item \eqref{eq:vi_start_collect_42}：
  若存在 $t'$ 使得 $P_{00}+I_{00}/k < \sum_{\tau=1}^{t'} d_\tau$，则仅靠初始库存无法满足到 $t'$ 的累积需求。
  在模型不允许欠交付（无 backlogging）且无外购原料的前提下，必须在 $t'$ 及之前产生至少一次原料“流入”，也就是至少一次收集访问（$z_{i\tau}=1$）。
  \item \eqref{eq:vi_totalcap_43}：
  任一可行路由方案在时期 $t$ 的总载重不超过 $KQ$，而左边正是主问题计划的总收集量，因此是必要条件。
  \item \eqref{eq:vi_bigitem_44}：
  若某条路径容量为 $Q$，则不可能同时装下两个需求都 $>Q/2$ 的客户，因此“大客户”的数量至多等于可用车辆数 $K$。
\end{itemize}

\paragraph{C.2 初始割（B.2--B.3）不会切掉最优解}
\begin{itemize}[leftmargin=2em]
  \item \eqref{eq:initcut_t1_49} 来自 DualCVRP(t) 的可行解：DualT1 的可行域是 DualCVRP(t) 的子集（仅额外要求 $w\ge0$），因此 $(w^*,u_0^*)$ 对 DualCVRP(t) 必可行。弱对偶给出 $\sum \lambda w^*+K u_0^*$ 是 CVRP(t) 的下界，从而 $\omega_t$ 作为 CVRP(t) 最优值必须不小于它。
  \item \eqref{eq:initcut_t2_412}：
  $T2$ 的对偶可行解在每个周期 $t$ 上的分量 $(w^*_{\cdot\cdot t},u^*_{0t})$ 满足与 DualCVRP(t) 同形的约束 \eqref{eq:t2_dual_constr_411}，故同样可作为每期 CVRP(t) 的对偶可行解并给出下界。将各期下界求和即可得到 \eqref{eq:initcut_t2_412}。
\end{itemize}

% ------------------------------------------------------------
% D. Implementation checklist
% ------------------------------------------------------------
\subsection{D. 代码实现 Checklist（建议按此顺序做）}

\begin{algorithm}[H]
\caption{把第4章加速策略接入 LBBD 的最小实现流程}
\DontPrintSemicolon
\KwIn{基础数据 + 第2.4节预计算参数 $\pi(i,t),g_{ivt}$（以及可选的 $e_{ivt}$）}
\KwIn{已实现的 LBBD 框架（BMP/BSP/列生成与割回调）}
\BlankLine
\textbf{(0) 主问题建模阶段：添加有效不等式}\;
计算 $t'$（若存在）并加入 \eqref{eq:vi_start_collect_42}\;
对每个 $t\in T$ 加入 \eqref{eq:vi_totalcap_43} 与 \eqref{eq:vi_bigitem_44}\;
\BlankLine
\textbf{(1) Benders 迭代前：生成初始割（可选，但推荐）}\;
\If{使用 $T1$ 初始割}{
对每个 $t\in T$ 解一次 DualT1（或解 $T1$ 并读出对偶变量），得到 $(w^*,u_0^*)$\;
向 BMP 加入 \eqref{eq:initcut_t1_49}\;
}
\If{使用 $T2$ 初始割}{
解一次全局 LP 松弛 $T2$ 并读出对偶可行解 $(w^*,u_0^*)$\;
向 BMP 加入 \eqref{eq:initcut_t2_412}\;
}
\BlankLine
\textbf{(2) 进入常规 LBBD 迭代：}\;
重复：解 BMP $\rightarrow$ 解各期 BSP(CVRP) $\rightarrow$ 生成并加入可行性/最优性割，直到收敛\;
\end{algorithm}

\end{document}
 或单独编译均可

\documentclass{article}
\usepackage{amsmath,amssymb,amsfonts,amsthm}
\usepackage{enumitem}
\usepackage{geometry}
\usepackage{mathrsfs}
\geometry{a4paper, margin=1in}
\usepackage{ctex} % 中文
\usepackage{hyperref}
\usepackage{float}
\usepackage[linesnumbered,ruled,vlined]{algorithm2e}

\newtheorem{remark}{注}
\newtheorem{definition}{定义}

\begin{document}

% ============================================================
%  Speed-up strategies for LBBD (Chapter 4)
%  1) Valid inequalities for BMP (master)
%  2) Warm-start initial cuts via dual-feasible solutions
% ============================================================

\section{LBBD 算法加速策略（第4章：有效不等式 + 初始割）}

\subsection{符号速查（只列本文件会用到的）}

\begin{itemize}[leftmargin=2em]
  \item 规划期：$T=\{1,\dots,l\}$。
  \item 收集点集合：$V_S$。在路由子问题中，客户集合写作 $N_c$（通常可取 $N_c=V_S$ 或“当期被访问的收集点”）。
  \item 车辆：可用车辆数上限 $K$，单车容量 $Q$。
  \item 主问题变量：
    \begin{itemize}[leftmargin=2em]
      \item $z_{it}\in\{0,1\}$：时期 $t$ 是否访问收集点 $i$。
      \item $\lambda_{ivt}\in\{0,1\}$：收集点 $i$ 的“上次访问时期为 $v$、本次访问时期为 $t$”的弧是否被选中（时间扩展网络/最短路式改写中的弧选择变量）。
      \item $\omega_t\ge 0$：时期 $t$ 的最优路径成本（由子问题给出，用 Benders 割约束它）。
    \end{itemize}
  \item 预计算参数（来自第2.4节改写模型；加速策略只\textbf{引用}它们，不重复推导）：
    \begin{itemize}[leftmargin=2em]
      \item $\pi(i,t)$：在“本次访问 $t$”条件下，\emph{可行的}最早上次访问时期下界（由收集点容量等逻辑得到）。
      \item $g_{ivt}$：若选择弧 $(v,t)$，即 $i$ 上次在 $v$ 被访问、当前在 $t$ 被访问，则时期 $t$ 从 $i$ 需要收走的\textbf{全量}（累计）收集量。
      \item （若实现 $T2$）$e_{ivt}$：与弧 $(v,t)$ 关联的收集点库存持有成本（或其它在第2.4节中被“弧化”的成本项）。
    \end{itemize}
\end{itemize}

\begin{remark}
本文件只覆盖第4章“加速方法”，默认你已经按 \texttt{lbbd\_all.tex} 中的说明实现了：
(1) BMP（主问题）中用 $\omega_t$ 代替显式路由变量；(2) BSP（子问题）用集合划分/列生成求解 CVRP(t) 并生成 Benders 割。
\end{remark}

% ------------------------------------------------------------
% A. Valid inequalities (add to BMP)
% ------------------------------------------------------------
\subsection{A. 主问题（BMP）有效不等式：直接加到主问题即可}

\subsubsection*{A.1 最小收集启动时间约束（式(4-1)--(4-2)）}

定义一个最早必须启动收集或生产的时期 $t'$：
\begin{equation}
t'=\arg\min_{1\le t\le l}\left\{\,P_{00}+\frac{I_{00}}{k}<\sum_{\tau=1}^{t} d_{\tau}\right\}.
\label{eq:tprime_def_41}
\end{equation}

若这样的 $t'$ 存在，则到 $t'$ 期末必须至少发生一次收集活动（否则没有新增原料来源，无法满足累积需求）：
\begin{equation}
\sum_{\tau=1}^{t'}\ \sum_{i\in V_S} z_{i\tau}\ \ge\ 1.
\label{eq:vi_start_collect_42}
\end{equation}

\paragraph{实现提示}
\begin{itemize}[leftmargin=2em]
  \item 这是一个\textbf{一次性}约束：只需在建模时计算 $t'$ 并加入一次 \eqref{eq:vi_start_collect_42}。
  \item 若不存在 $t'$（即初始库存 $P_{00}+I_{00}/k$ 足以覆盖全周期累积需求），则不要添加该约束。
\end{itemize}

\subsubsection*{A.2 周期车辆总容量约束（式(4-3)）}

时期 $t$ 的总计划收集量可由 $\lambda$ 的选择唯一确定（全量收集）：
\[
\text{Pickup}_t \;=\; \sum_{i\in V_S}\ \sum_{v=\pi(i,t)}^{t-1} g_{ivt}\,\lambda_{ivt}.
\]
由于每期最多使用 $K$ 辆容量为 $Q$ 的车，必要条件为：
\begin{equation}
\sum_{i\in V_S}\ \sum_{v=\pi(i,t)}^{t-1} g_{ivt}\,\lambda_{ivt}\ \le\ K\cdot Q,\qquad \forall t\in T.
\label{eq:vi_totalcap_43}
\end{equation}

\paragraph{实现提示}
\begin{itemize}[leftmargin=2em]
  \item 这是\textbf{按期}约束：对所有 $t\in T$ 添加 \eqref{eq:vi_totalcap_43}。
  \item 该不等式是可行性的\textbf{必要条件}（不是充分条件），但能显著收紧主问题的 LP 松弛。
\end{itemize}

\subsubsection*{A.3 大收集量客户约束（式(4-4)）}

若某个收集点在时期 $t$ 的收集量超过半车容量（$g_{ivt}>Q/2$），则同一条路径不可能再服务另一个同样“大”的收集点。
因此，在时期 $t$，满足 $g_{ivt}>Q/2$ 的被选中弧数量不能超过 $K$：
\begin{equation}
\sum_{i\in V_S}\ \sum_{\substack{v=\pi(i,t)\\ g_{ivt}>Q/2}}^{t-1}\ \lambda_{ivt}\ \le\ K,\qquad \forall t\in T.
\label{eq:vi_bigitem_44}
\end{equation}

\paragraph{实现提示}
\begin{itemize}[leftmargin=2em]
  \item 这同样是\textbf{按期}约束：对每个 $t$ 建一个线性约束，内部只枚举满足 $g_{ivt}>Q/2$ 的 $(i,v)$。
  \item 若存在 $g_{ivt}>Q$，则该弧本身会导致 CVRP(t) 不可行，建议在预处理阶段直接\textbf{删除}该弧或在主问题中禁用它（比加 \eqref{eq:vi_bigitem_44} 更“直接”）。
\end{itemize}

% ------------------------------------------------------------
% B. Initial cuts (warm start)
% ------------------------------------------------------------
\subsection{B. 主问题（BMP）初始割：Benders 迭代前的“预热”}

动机：LBBD 初期，$\omega_t$ 在主问题里可能被压到非常小（甚至 0），导致第一次/前几次迭代的主问题解偏离真实可行解，从而迭代次数增多。初始割为 $\omega_t$ 提供一个\textbf{有效下界}，加速收敛。

\subsubsection*{B.1 访问决策未知的 CVRP(t) 集合划分模型（式(4-5)）}

设 $\hat{\mathcal R}_t$ 为时期 $t$ 的所有可行路径集合（通过列生成隐式表示），路径 $r$ 的成本为 $b_r$，其对“客户-需求情景”$(i,v)$ 的覆盖系数为 $a_{ivr}\in\{0,1\}$。
则 CVRP(t) 可写为：
\begin{equation}
\begin{aligned}
\text{CVRP}(t)\;\; \omega_t=\min\ &\sum_{r\in\hat{\mathcal R}_t} b_r\,\xi_r\\
\text{s.t.}\ &\sum_{r\in\hat{\mathcal R}_t} a_{ivr}\,\xi_r=\lambda_{ivt},\quad \forall i\in N_c,\ v=\pi(i,t),\dots,t-1,\\
&\sum_{r\in\hat{\mathcal R}_t}\xi_r\le K,\\
&\xi_r\in\{0,1\},\quad \forall r\in\hat{\mathcal R}_t.
\end{aligned}
\label{eq:cvrp_sp_45}
\end{equation}

其 LP 松弛对偶为（式(4-6)）：
\begin{equation}
\begin{aligned}
\text{DualCVRP}(t)\ \max\ &\sum_{i\in N_c}\sum_{v=\pi(i,t)}^{t-1} \lambda_{ivt}\,w_{iv}+K\,u_0\\
\text{s.t.}\ &\sum_{i\in N_c}\sum_{v=\pi(i,t)}^{t-1} a_{ivr}\,w_{iv}+u_0\le b_r,\quad \forall r\in\hat{\mathcal R}_t,\\
&w_{iv}\in\mathbb R,\quad \forall i\in N_c,\ v=\pi(i,t),\dots,t-1,\\
&u_0\le 0.
\end{aligned}
\label{eq:dual_cvrp_46}
\end{equation}

\begin{remark}
对任意 DualCVRP(t) 的可行解 $(w,u_0)$，其目标值都是 CVRP(t) 的\textbf{下界}（弱对偶）。
因此在 BMP 中添加 $\omega_t\ge \sum \lambda\,w+K u_0$ 形式的不等式一定是\textbf{安全的初始割}。
\end{remark}

\subsubsection*{B.2 初始割I：基于单周期松弛模型 $T1$ 的对偶（式(4-7)--(4-9)）}

为快速构造 DualCVRP(t) 的一个可行解，论文给出一个更容易（在实践中更“稳”）的松弛：
\begin{equation}
\begin{aligned}
(T1)\ \min\ &\sum_{r\in\hat{\mathcal R}_t} b_r\,\xi_r\\
\text{s.t.}\ &\sum_{r\in\hat{\mathcal R}_t} a_{ivr}\,\xi_r\ge 1,\quad \forall i\in N_c,\ v=\pi(i,t),\dots,t-1,\\
&\sum_{r\in\hat{\mathcal R}_t}\xi_r\le K,\\
&\xi_r\ge 0,\quad \forall r\in\hat{\mathcal R}_t.
\end{aligned}
\label{eq:t1_primal_47}
\end{equation}

其对偶（式(4-8)）为：
\begin{equation}
\begin{aligned}
(\text{Dual }T1)\ \max\ &\sum_{i\in N_c}\sum_{v=\pi(i,t)}^{t-1} w_{iv}+K u_0\\
\text{s.t.}\ &\sum_{i\in N_c}\sum_{v=\pi(i,t)}^{t-1} a_{ivr}\,w_{iv}+u_0\le b_r,\quad \forall r\in\hat{\mathcal R}_t,\\
&w_{iv}\ge 0,\quad \forall i\in N_c,\ v=\pi(i,t),\dots,t-1,\\
&u_0\le 0.
\end{aligned}
\label{eq:dual_t1_48}
\end{equation}

由于 \eqref{eq:dual_t1_48} 与 \eqref{eq:dual_cvrp_46} 的约束相同但更严格（$w_{iv}\ge 0$），所以 DualT1 的任意可行解也是 DualCVRP(t) 的可行解。
令 $(w^*,u_0^*)$ 为 DualT1 的任意一个可行解（通常取其最优解），则可加入初始割（式(4-9)）：
\begin{equation}
\omega_t\ \ge\ \sum_{i\in N_c}\sum_{v=\pi(i,t)}^{t-1} w_{iv}^*\,\lambda_{ivt}\ +\ K u_0^*.
\label{eq:initcut_t1_49}
\end{equation}

\paragraph{实现提示（怎样在代码里“拿到”$(w^*,u_0^*)$）}
\begin{itemize}[leftmargin=2em]
  \item 你可以直接解 \eqref{eq:dual_t1_48}（对偶）或 \eqref{eq:t1_primal_47}（原始）得到对偶变量。
  \item 由于 $\hat{\mathcal R}_t$ 很大，通常仍需要列生成：
    \begin{itemize}[leftmargin=2em]
      \item 给定 $(w,u_0)$，路径 $r$ 的对偶约束违反/定价等价于检查\;
      $\min_{r}\ \bigl(b_r-\sum_{i,v}a_{ivr}w_{iv}-u_0\bigr)$ 是否为负。
      \item 这与 CVRP 的“最短路/定价子问题”同源，可复用你已有的定价器。
    \end{itemize}
  \item 为了让系数更“有用”，论文建议采用 $\ge 1$（从而 $w_{iv}\ge 0$）而不是等式 $=1$（否则 $w_{iv}$ 无符号，容易给出较弱/不稳定的 cut 系数）。
\end{itemize}

\begin{remark}[$T1$ 初始割的全局有效性与强度权衡]
$T1$ 将覆盖约束从 $=\lambda_{ivt}$ 改为 $\ge 1$，使得 $T1$ 的求解\textbf{不依赖 $\lambda$ 的具体取值}——对所有情景节点 $(i,v)$ 均要求至少覆盖一次。这一特性带来两个直接后果：
\begin{itemize}
  \item \textbf{全局有效性}：割系数 $w^*_{iv}$ 对所有可能的 $\lambda$ 配置均成立，割 \eqref{eq:initcut_t1_49} 作为全局有效不等式可安全加入 BMP；
  \item \textbf{强度偏弱}：由于 $T1$ 忽略了访问序列约束（哪些弧实际被选中），其对偶解未能利用 $\lambda$ 结构，给出的下界通常比基于完整 $\lambda$ 信息的 $T2$ 割更松。
\end{itemize}
因此，$T1$ 割适合作为快速预热手段；若计算资源允许，建议进一步叠加 $T2$ 割以获得更紧的初始下界。
\end{remark}

\subsubsection*{B.3 初始割II：基于全局松弛模型 $T2$ 的对偶（式(4-10)--(4-12)）}

思想：$T1$ 忽略了很多来自“生产-库存-访问序列”的全局约束信息，得到的对偶解可能偏松；
因此构造一个更接近完整模型的 LP 松弛 $T2$，从其对偶可行解中抽取每期的 $(w^*_{ivt},u^*_{0t})$。

论文给出的 $T2$（式(4-10)）是“将第2.4节重构后的完整模型做 LP 松弛，并将路由部分用集合划分表示”：
\begin{equation}
\begin{aligned}
(T2)\ \min\ &\sum_{t\in T}\left(\sum_{r\in\hat{\mathcal R}_t} b_r\,\xi_{rt} + u p_t + f y_t + h_0 I_{0t}+h_p P_{0t}\right)
+\sum_{t\in T'}\sum_{i\in V_S}\sum_{v=\pi(i,t)}^{t-1} e_{ivt}\lambda_{ivt}\\
\text{s.t.}\ &\sum_{r\in\hat{\mathcal R}_t} a_{ivrt}\,\xi_{rt}=\lambda_{ivt},\quad \forall i\in N_c,\ t\in T,\ v=\pi(i,t),\dots,t-1,\\
&\sum_{r\in\hat{\mathcal R}_t}\xi_{rt}\le K,\quad \forall t\in T,\\
&0\le y_t\le 1,\ \forall t\in T,\qquad 0\le \lambda_{ivt}\le 1,\ \forall(i,v,t),\\
&\xi_{rt}\ge 0,\ \forall r\in\hat{\mathcal R}_t,\ t\in T,\\
&\text{并满足（重构模型中的其余线性约束），例如论文中的}\ (2\text{-}37)\sim(2\text{-}45),(2\text{-}60)\sim(2\text{-}62).
\end{aligned}
\label{eq:t2_lp_410}
\end{equation}

其对偶中，与 $\xi_{rt}$ 相关的对偶约束为（式(4-11)）：
\begin{equation}
\sum_{i\in N_c}\sum_{v=\pi(i,t)}^{t-1} a_{ivrt}\,w_{ivt}+u_{0t}\le b_r,\qquad \forall r\in\hat{\mathcal R}_t,\ t\in T.
\label{eq:t2_dual_constr_411}
\end{equation}

记 $(w^*,u_0^*)$ 为 $T2$ 的任意一个对偶可行解（通常取最优对偶解）。
论文给出的更紧初始割（式(4-12)）为：
\begin{equation}
\sum_{t\in T}\omega_t\ \ge\ \sum_{t\in T}\left(\sum_{i\in N_c}\sum_{v=\pi(i,t)}^{t-1} w_{ivt}^*\,\lambda_{ivt} + K u_{0t}^*\right).
\label{eq:initcut_t2_412}
\end{equation}

\paragraph{实现提示}
\begin{itemize}[leftmargin=2em]
  \item 与 $T1$ 相比，$T2$ 规模更大、求解更慢，但系数通常更紧。
  \item 论文形式写成“总和 cut” \eqref{eq:initcut_t2_412}（一次加一条），实现最简单：在开始迭代前解一次 $T2$ 的 LP，取对偶变量加一条不等式即可。
  \item 若你愿意，也可以把 \eqref{eq:initcut_t2_412} 拆成每期一条 $\omega_t\ge \sum w^*_{ivt}\lambda_{ivt}+K u^*_{0t}$（更紧），但这不属于论文原式，需自行做数值稳定性测试。
\end{itemize}

% ------------------------------------------------------------
% C. Correctness checks
% ------------------------------------------------------------
\subsection{C. 正确性检查（这些加速策略为什么“安全”）}

\paragraph{C.1 有效不等式（A.1--A.3）不会切掉最优整数解}
\begin{itemize}[leftmargin=2em]
  \item \eqref{eq:vi_start_collect_42}：
  若存在 $t'$ 使得 $P_{00}+I_{00}/k < \sum_{\tau=1}^{t'} d_\tau$，则仅靠初始库存无法满足到 $t'$ 的累积需求。
  在模型不允许欠交付（无 backlogging）且无外购原料的前提下，必须在 $t'$ 及之前产生至少一次原料“流入”，也就是至少一次收集访问（$z_{i\tau}=1$）。
  \item \eqref{eq:vi_totalcap_43}：
  任一可行路由方案在时期 $t$ 的总载重不超过 $KQ$，而左边正是主问题计划的总收集量，因此是必要条件。
  \item \eqref{eq:vi_bigitem_44}：
  若某条路径容量为 $Q$，则不可能同时装下两个需求都 $>Q/2$ 的客户，因此“大客户”的数量至多等于可用车辆数 $K$。
\end{itemize}

\paragraph{C.2 初始割（B.2--B.3）不会切掉最优解}
\begin{itemize}[leftmargin=2em]
  \item \eqref{eq:initcut_t1_49} 来自 DualCVRP(t) 的可行解：DualT1 的可行域是 DualCVRP(t) 的子集（仅额外要求 $w\ge0$），因此 $(w^*,u_0^*)$ 对 DualCVRP(t) 必可行。弱对偶给出 $\sum \lambda w^*+K u_0^*$ 是 CVRP(t) 的下界，从而 $\omega_t$ 作为 CVRP(t) 最优值必须不小于它。
  \item \eqref{eq:initcut_t2_412}：
  $T2$ 的对偶可行解在每个周期 $t$ 上的分量 $(w^*_{\cdot\cdot t},u^*_{0t})$ 满足与 DualCVRP(t) 同形的约束 \eqref{eq:t2_dual_constr_411}，故同样可作为每期 CVRP(t) 的对偶可行解并给出下界。将各期下界求和即可得到 \eqref{eq:initcut_t2_412}。
\end{itemize}

% ------------------------------------------------------------
% D. Implementation checklist
% ------------------------------------------------------------
\subsection{D. 代码实现 Checklist（建议按此顺序做）}

\begin{algorithm}[H]
\caption{把第4章加速策略接入 LBBD 的最小实现流程}
\DontPrintSemicolon
\KwIn{基础数据 + 第2.4节预计算参数 $\pi(i,t),g_{ivt}$（以及可选的 $e_{ivt}$）}
\KwIn{已实现的 LBBD 框架（BMP/BSP/列生成与割回调）}
\BlankLine
\textbf{(0) 主问题建模阶段：添加有效不等式}\;
计算 $t'$（若存在）并加入 \eqref{eq:vi_start_collect_42}\;
对每个 $t\in T$ 加入 \eqref{eq:vi_totalcap_43} 与 \eqref{eq:vi_bigitem_44}\;
\BlankLine
\textbf{(1) Benders 迭代前：生成初始割（可选，但推荐）}\;
\If{使用 $T1$ 初始割}{
对每个 $t\in T$ 解一次 DualT1（或解 $T1$ 并读出对偶变量），得到 $(w^*,u_0^*)$\;
向 BMP 加入 \eqref{eq:initcut_t1_49}\;
}
\If{使用 $T2$ 初始割}{
解一次全局 LP 松弛 $T2$ 并读出对偶可行解 $(w^*,u_0^*)$\;
向 BMP 加入 \eqref{eq:initcut_t2_412}\;
}
\BlankLine
\textbf{(2) 进入常规 LBBD 迭代：}\;
重复：解 BMP $\rightarrow$ 解各期 BSP(CVRP) $\rightarrow$ 生成并加入可行性/最优性割，直到收敛\;
\end{algorithm}

\end{document}
